% MINT method figure (IEEE two-column, full width, drawn at final size: no scaling).
% Needs: \usepackage{tikz,xcolor,amsmath,amssymb}; \usetikzlibrary{arrows.meta,positioning,fit,calc}
% (all loaded by main.tex). Okabe-Ito palette (colour-blind safe); every colour cue is doubled by a
% line style or a text label, so the figure also reads in grayscale.
% Layout: (a) the five stages of MINT for one receiver frame, 5 boxes of 2.05 cm at a 2.35 cm pitch;
% (b) a top-down example of stage 3. The decision rules themselves are in the text and Algorithm 1,
% not repeated here. Total width 17.7 cm (IEEE \textwidth = 18.1 cm).
\begin{figure*}[!t]
\centering
\definecolor{mintRecv}{HTML}{0072B2}   % receiver / trusted evidence (blue)
\definecolor{mintMal}{HTML}{D55E00}    % untrusted sender, TRAV (vermillion)
\definecolor{mintSup}{HTML}{009E73}    % SUP (bluish green)
\newcommand{\mintStep}[1]{{\setlength{\fboxsep}{0.8pt}\colorbox{black!85}{\textcolor{white}{\bfseries\,#1\,}}}}
\begin{tikzpicture}[
  font=\scriptsize,
  blk/.style={draw, rounded corners=1.5pt, align=left, anchor=north west, inner sep=2.5pt,
              line width=0.5pt, text width=1.9cm, minimum height=3.05cm},
  stage/.style={anchor=north west, align=left, inner sep=2.5pt, text width=1.9cm},
  arr/.style={-{Latex[length=1.5mm]}, line width=0.6pt},
  trust/.style={-{Latex[length=1.5mm]}, line width=0.9pt, mintRecv}]

%---------------------------------------------------------------- (a) five stages
\node[anchor=south west, inner sep=0pt] at (0,0.62) {\textbf{(a) \mint{} for one receiver frame}};
\node[blk, draw=mintRecv, fill=mintRecv!7, line width=0.8pt, minimum height=1.2cm] (rx) at (0,0)
  {\mintStep{1} \textbf{Inputs}\\[1pt] receiver: scan $P_0$,\\ features $f_0$ (\textbf{trusted})};
\node[blk, draw=mintMal, fill=mintMal!5, dashed, line width=0.8pt, minimum height=1.7cm] (tx) at (0,-1.35)
  {senders $i\ge1$:\\ features $f_i$\\ (\textbf{untrusted})};
\node[stage] (s2) at (2.35,0)
  {\mintStep{2} \textbf{Counterfactual attribution}\\[2pt] decode with all agents ($Y^{+}$) and without
   each sender ($Y^{-}(i)$); differences are ADD, REMOVE or SHIFT effects of $i$ in tile $u$};
\node[stage] (s3) at (4.7,0)
  {\mintStep{3} \textbf{Receiver-ray verification}\\[2pt] count the receiver's own rays that end in
   (\textcolor{mintSup!80!black}{\textbf{SUP}}) or pass through (\textcolor{mintMal}{\textbf{TRAV}}) each
   effect box; too few rays: UNKNOWN};
\node[stage] (s4) at (7.05,0)
  {\mintStep{4} \textbf{Gating decision}\\[2pt] switch sender $i$ off in tile $u$ only if the rays
   resolvably contradict an effect and confirm none; UNKNOWN keeps};
\node[stage] (s5) at (9.4,0)
  {\mintStep{5} \textbf{Filtered detections}\\[2pt] each tile takes the decode without its switched-off
   senders; NMS gives $Y_{\mathrm{MINT}}$};

% stage frames: drawn explicitly so that all four have the same height
\foreach \x in {2.35,4.7,7.05} {\draw[black!75, rounded corners=1.5pt, line width=0.5pt] (\x,0) rectangle ++(2.05,-3.05);}
\draw[black!75, rounded corners=1.5pt, line width=0.9pt] (9.4,0) rectangle ++(2.05,-3.05);
\draw[arr] (rx.east) -- (rx.east -| 2.35,0);
\draw[arr, mintMal, dashed] (tx.east) -- (tx.east -| 2.35,0);
\foreach \x in {4.4,6.75,9.05} {\draw[arr] (\x,-1.3) -- ++(0.3,0);}
% trusted evidence: the receiver's own scan reaches stage 3 directly, never via fusion
\draw[trust] ([xshift=0.5cm]rx.north) -- ++(0,0.3) -| (5.725,0)
  node[pos=0.25, above, text=mintRecv, inner sep=1pt] {own scan $P_0$, never fused};

%---------------------------------------------------------------- (b) top-down example of stage 3
\draw[black!50, fill=black!4, rounded corners=1.5pt, line width=0.5pt] (11.75,0.95) rectangle (17.7,-3.05);
\node[anchor=north west, inner sep=2.5pt] at (11.75,0.95) {\textbf{(b) Stage 3, seen from above}};
\begin{scope}[shift={(12.45,-1.6)}, scale=0.85, every node/.style={scale=0.85}]
  \draw[black!20] (0.9,0) grid[step=1.25] (4.65,2.5);
  \draw[black!80, line width=0.9pt] (2.15,1.25) rectangle (3.4,2.5);
  \node[anchor=north east, font=\footnotesize] at (3.4,2.5) {$u$};
  \node[anchor=south east, font=\footnotesize] at (4.65,0) {$u'$};
  \coordinate (S) at (0.45,1.25);
  \draw[mintRecv, fill=mintRecv!30, line width=0.6pt] (0.12,1.08) rectangle (0.78,1.42);
  \fill[mintRecv] (S) circle (1.5pt);
  \node[anchor=north, text=mintRecv, font=\footnotesize] at (0.45,1.05) {receiver};
  \begin{scope}
    \clip (0.9,0) rectangle (4.65,2.5);
    \fill[black!14] (2.4,0.92) -- ($(S)!3!(2.4,0.92)$) -- ($(S)!3!(2.4,0.55)$) -- (2.4,0.55) -- cycle;
  \end{scope}
  \draw[black!60, densely dotted, line width=0.7pt] (3.7,0.1) rectangle (4.3,0.38);
  \node[font=\footnotesize] at (4.0,0.24) {?};
  \draw[mintMal, dashed, line width=0.8pt] (2.35,1.6) rectangle (3.05,1.95);
  \node[text=mintMal, font=\footnotesize, inner sep=1pt] at (2.6,2.24) {fake};
  \foreach \y in {1.64,1.72,1.80,1.88} {
    \draw[mintMal, densely dashed, line width=0.5pt] (S) -- ($(S)!1.75!(2.35,\y)$);
  }
  \draw[black, fill=black!40, line width=0.6pt] (2.4,0.55) rectangle (3.1,0.92);
  \node[anchor=north, font=\footnotesize, inner sep=1pt] at (2.75,0.55) {car};
  \foreach \y in {0.61,0.69,0.77,0.85} {
    \draw[mintSup, line width=0.6pt] (S) -- (2.4,\y);
  }
\end{scope}
\node[anchor=north west, align=left, text width=5.65cm, inner sep=0pt] at (11.9,-1.75)
  {\textcolor{mintMal}{\textbf{TRAV}} (dashed): rays cross the fake box and end behind it, so the ADD is
   contradicted. \textcolor{mintSup!80!black}{\textbf{SUP}} (solid): rays end on the real car, which
   confirms it. Shaded: the receiver cannot see here ($N{<}3$), so the evidence is UNKNOWN and the sender is
   kept.};
\end{tikzpicture}
\caption{Overview of \mint. (a) For each receiver frame, the frozen detector decodes once with all agents and
once without each sender. The differences are attributed to a sender and a 6.4\,m tile. Each difference is
checked against the receiver's own LiDAR rays, which never pass through fusion. A sender is switched off only in
the tiles where the rays clearly contradict it, and the output is assembled tile by tile. (b) A fake car added by
a sender is crossed by the receiver's rays, while a real car stops them; a region the receiver cannot see gives
no evidence and is left unchanged.}
\label{fig:mintmethod}
\end{figure*}
